\chapter{第二讲\quad 半线性波动方程的应用}

我们关注重要几何系统(如波动映射、Yang--Mills 与 Maxwell--Klein--Gordon 方程)的柯西问题. 作为第一近似, 所有这些方程都可写成形式
\begin{equation}
\Box\phi=\mathcal{N}(\phi),\end{equation}
其中 $\Box = -\partial_t^2 + \Delta$ 是 $\mathbb{R} \times \mathbb{R}^n$ 上的标准波动算子, $\Delta = \sum_1^n \partial_j^2$ 是 $\mathbb{R}^n$ 上的 Laplacian, $\phi = \phi(t, x)$ 取值于某 $N \geq 1$ 的 $\mathbb{R}^N$, 而 $\mathcal{N}$ 是关于时间变量局部的非线性算子(即 $\mathcal{N}(\phi)(t, \cdot)$ 只依赖于 $\phi(t', \cdot)$, 其中 $t'$ 在 $t$ 的任意邻域中).
我们在初值超平面 $t = 0$ 上给定柯西数据:
\begin{equation}
(\phi,\partial_{t}\phi)_{t=0}=(f,g)\in H^{s}\times H^{s-1},\end{equation}
其中 $H^s = \{f : (I - \Delta)^{s/2} f \in L^2\}$. 为简便记 $\phi[0] \in H^s$.
我们事实上集中关注以下类型的系统:
\section*{1. 波动映射(Wave Maps)}

\begin{equation}
\Box\phi^{I}+\sum_{J,K}\Gamma_{J K}^{I}(\phi)Q_{0}(\phi^{J},\phi^{K})=0.\end{equation}

这里 $\phi^{I}$ 表示 $\phi$ 的第 $I$ 个分量函数, $\Gamma_{JK}^{I}$ 是从 $\mathbb{R}^{N}$ 到 $\mathbb{R}$ 的光滑函数, $Q_{0}$ 为零形式
\[
Q_{0}(\phi,\psi)=\sum_{\mu=0}^{n}\partial_{\mu}\phi\partial^{\mu}\psi=-\partial_{t}\phi\partial_{t}\psi+\sum_{j=1}^{n}\partial_{j}\phi\partial_{j}\psi.
\]

\section*{2. Yang--Mills 型}

\begin{equation}
\Box\phi=D^{-1}Q(\phi,\phi)+Q(D^{-1}\phi,\phi),\end{equation}
其中 $D^\alpha = (-\Delta)^{\alpha/2}$, $Q$ 表示如下类型的任一双线性算子: 给定向量值函数 $\phi$ 与 $\psi$, $Q(\phi, \psi)$ 的第 $I$ 个分量函数是对所有 $1 \leq i < j \leq n$ 与所有 $J, K$ 的 $Q_{ij}(\phi^J, \psi^K)$ 的常实系数线性组合, 其中 $Q_{ij}$ 为零形式
\[
Q_{i j}(\phi,\psi)=\partial_{i}\phi\partial_{j}\psi-\partial_{i}\psi\partial_{j}\phi.
\]
((YM) 右端的两个 $Q$ 可能代表两个不同的此类算子. )

\section*{3. Maxwell--Klein--Gordon 型}

\begin{equation}
\left\{\begin{aligned}\Box\phi&=D^{-1}Q(\psi,\psi),\\ \Box\psi&=Q(D^{-1}\phi,\psi),\end{aligned}\right.\end{equation}
其中 $\phi = (\phi^1, \ldots, \phi^{N_1})$, $\psi = (\psi^1, \ldots, \psi^{N_2})$, $N = N_1 + N_2$, 而 $Q$ 意义同前. 因此 (MKG) 是 (YM) 的特例.
下面的定理概括了当今已知的主要适定性结果.
\begin{theorem*}[主定理]
\begin{enumerate}[label=\roman*.]
\item (\cite{klai35,klai44}. )若 $n \geq 2$ 且 $s > s_c = \frac{n}{2}$, 则 (WM) 对 $H^s(\mathbb{R}^n)$ 中的初值局部适定.
\item (\cite{klai72}、\cite{klai73}. )对 $n \geq 2$, 系统 (WM) 在 Besov 空间 $B_{2}^{\frac{n}{2},1}(\mathbb{R}^{n})$ 中对小数据局部适定.
\item (\cite{klai70}、\cite{klai41}、\cite{klai71}. )若 $n \geq 5$ 且目标流形承认有界可平行化结构, 则波动映射系统在 $H^{\frac{n}{2}}(\mathbb{R}^{n})$ 中对小数据"弱适定". 若目标流形为标准球面 $\mathbf{S}^{k}$, 则对 $n \geq 2$ 同样成立.
\item (\cite{klai33}、\cite{klai34}、\cite{klai19}. )若 $n = 3$, 则完整的 MKG 与 YM 系统在能量范数 $H^{2}(\mathbb{R}^{n})$ 下对大初值整体适定. 对临界幂非线性的 Yang--Mills--Higgs 系统结果亦成立.
\item (\cite{klai13}、\cite{klai21}. )若 $n = 3$, 则完整 MKG 系统在 $s > \frac{3}{4}$ 时局部适定, 在 $s > \frac{7}{8}$ 时整体适定.
\item (\cite{klai39}、\cite{klai47}. )若 $n = 4$, 则约化 MKG 与 YM 型系统在 $s > s_c = 1$ 时局部适定.\end{enumerate}
\end{theorem*}

\section*{2.1 主定理的动机}

考虑系统
\begin{equation}
\Box u=F(u,\partial u),\end{equation}
其中 $u : \mathbb{R}^{1+n} \to \mathbb{R}^N$, $\partial u = (\partial_t u, \partial_1 u, \ldots, \partial_n u)$, $F$ 为满足 $F(0) = 0$ 的光滑 $\mathbb{R}^N$ 值函数. 对这个方程有如下标准存在唯一性结果:
\setcounter{mtthm}{1}% 2.1 为节标题, 下一定理环境显示 2.2

\begin{theorem}\label{thm:2.2}
方程 (22) 对 $H^s \times H^{s-1}(\mathbb{R}^n)$ 中的初值局部适定, 只要 $s > \frac{n}{2} + 1$. 即对任意初值 $u[0] \in H^s(\mathbb{R}^n)$($s > \frac{n}{2} + 1$), 存在 $T > 0$ 与唯一解
\begin{equation}
u\in C^{0}([0,T],H^{s}(\mathbb{R}^{n}))\cap C^{1}([0,T],H^{s-1}(\mathbb{R}^{n})).\end{equation}
此外解连续(事实上是解析地)依赖于数据.\end{theorem}

\begin{remark}\label{rem:2.3}
这一非常经典结果的证明只基于简单的能量估计、嵌入 $L^{\infty} \subset H^{s}$ 的 Sobolev 估计, 以及一些简单的插值不等式. 而且证明可无大困难地推广到极一般的演化方程类, 包括拟线性波动方程组.\end{remark}

\begin{remark}\label{rem:2.4}
就初值的正则性假设而言, 该结果远非最优. 最优适定性的问题至关重要, 并在过去 10 年中产生了大量工作.\end{remark}

为在我们的例子(波动映射、Maxwell--Klein--Gordon 与 Yang--Mills 方程)背景下更好地理解最优局部适定性问题, 我们需要定义临界适定性(以下缩写为 WP)指数 $s_c$. 我们所有方程都有与之相伴的自然伸缩, 而 $s_c$ 是使初数据的 $\dot{H}^s \times \dot{H}^{s-1}$ 范数在此伸缩下不变的 $s$ 的唯一值. 例如, 若 $u$ 满足 (WM), 则对任意 $\lambda > 0$,
\[
u_{\lambda}(t,x)=u(\lambda t,\lambda x)
\]
也满足. 由于 $\|u_{\lambda}(t)\|_{H^{s}} = \lambda^{\frac{n}{2}-s}\|u(\lambda t)\|_{H^{s}}$, (WM) 的临界 WP 指数为 $s_{c} = \frac{n}{2}$.
同一原理适用于 (MKG)、(YM). 事实上它们临界 WP 指数均为 $s_{c} = \frac{n-2}{2}$.
据此我们给出如下猜想, 见 \cite{klai29}:
\begin{theorem*}[一般适定性猜想]
\begin{enumerate}[label=1.]
\item 对所有基本场论, 初值问题对 $H^{s}$($s > s_c$)中的初值局部适定.
\item 基本场论对所有具有小 $H^{s_c} \times H^{s_c-1}$ 范数的初数据弱$^{5}$ 整体适定.
\item 基本场论对 $H^{s} \times H^{s-1}$($s < s_c$)中的初值不适定.\end{enumerate}
\end{theorem*}

为证明 $s > s_c$ 时的局部适定性, 我们在适当的 Banach 空间中用 Picard 迭代. 考虑柯西问题
\begin{equation}
\Box u=\mathcal{N}(u),\qquad(u,\partial_{t}u)_{t=0}=(f,g).\end{equation}
第 $0$ 次迭代 $u_{0}$ 为解的无齐次部分:
\[
\Box u_{0}=0,\qquad(u,\partial_{t}u)_{t=0}=(f,g).
\]
后续迭代归纳地由
\[
u_{j+1}=u_{0}+\Box^{-1}\mathcal{N}(u_{j})\qquad (j \geq 0)
\]
给出, 其中 $\Box^{-1}$ 是把任一充分正则的 $F$ 映为方程 $\Box v = F$ 满足 $(v, \partial_t v)_{t=0} = 0$ 的解 $v$ 的算子.
若要通过迭代证明 $\Box u = \mathcal{N}(u)$ 在 $H^s \times H^{s-1}$ 中初值的局部解存在, 必须能证明迭代保持在数据空间中:
\[
f\in H^{s},\quad g\in H^{s-1}\implies u_{j}(t)\in H^{s},\quad\partial_{t}u_{j}(t)\in H^{s-1}
\]
对所有 $j \geq 0$ 与某区间 $(0,T)$ 中所有 $t$ 成立. 对 $j=0$, (23) 显然; 但 $j=1$ 的情形已提供有价值的见解.
\setcounter{mtthm}{4}% 2.3、2.4 为注, 下一定义环境显示 2.5

\begin{definition}\label{def:2.5}
若 (23) 对 $j = 1$ 与所有 $(f, g) \in H^s \times H^{s-1}$ 成立, 我们称第一迭代对 $H^s$ 中的初值 WP.\end{definition}

\begin{example}\label{ex:2.1}
考虑模型问题
\begin{equation}
\Box u=(\partial_{t}u)^{2},\end{equation}
其中 $u$ 为实值. 该方程与 (WM) 有相同的伸缩性质, 故 WP 指数为 $s_c = \frac{n}{2}$. 我们想求使第一迭代 $u_1$ 对 $H^s$ 中初值 WP 的 $s$ 的集合的下界. 一个涉及 Duhamel 原理的简单计算表明, 这归结为证明形如
\[
\int_{\mathbb{R}^{n}\times\mathbb{R}^{n}}K(\xi,\eta)f(\xi)g(\eta)h(\xi+\eta)\,d\xi d\eta\lesssim\|f\|_{L^{2}}\|g\|_{L^{2}}\|h\|_{L^{2}}
\]
的估计, 对所有 $f, g, h \in L^{2}(\mathbb{R}^{n})$ 成立, 其中
\[
K\left(\xi,\eta\right)=\frac{\langle\xi+\eta\rangle^{s-1}}{\langle\xi\rangle^{s-1}\langle\eta\rangle^{s-1}\left(1+\Delta_{\pm}(\xi,\eta)\right)},
\]
\[
\Delta_{+}=|\xi|+|\eta|-|\xi+\eta|,\qquad\Delta_{-}=|\xi+\eta|-|\xi|-|\eta|.
\]
这里记号 $\langle\cdot\rangle=1+|\cdot|$.
关于形如
\[
\int_{\mathbb{R}^{n}\times\mathbb{R}^{n}}\frac{f(\xi)g(\eta)h(\xi+\eta)}{\langle\xi\rangle^{a}\langle\eta\rangle^{b}\bigl(1+\Delta_{\pm}(\xi,\eta)\bigr)^{c}}\,d\xi d\eta\lesssim\|f\|_{L^{2}}\|g\|_{L^{2}}\|h\|_{L^{2}}
\]
的积分估计, 可得如下结果, 其中 $\Delta_{\pm}$ 由 (26) 给出.\end{example}

\begin{proposition}\label{prop:2.6}
设 $a, b, c \geq 0$. 则 (27) 成立当且仅当 $a + b + c > \frac{n}{2}$ 且 $c < \frac{n-1}{4}$.\end{proposition}

\begin{remark}\label{rem:2.7}
应当指出, 基于简单的集中型反例, 当 $a + b + c < \frac{n}{2}$ 或 $c \geq \frac{n-1}{4}$ 时估计失败.\end{remark}

把\propref{prop:2.6} 用于核 (25) 需要
\[
s-1+\min\left(1,\frac{n-1}{4}\right)>\frac{n}{2},\qquad \text{即}\ s>\max\left(\frac{n}{2},\frac{n+5}{4}\right).
\]
因此, 对 $n = 3$ 维的模型方程 $\Box u = (\partial_t u)^2$, 上例表明第一迭代对 $H^s$ 中初值在 $s > 2 = s_c + \frac{1}{2}$ 时 WP; 事实上可证当 $s \leq 2$ 时这不成立. 这应与 Lindblad \cite{klai49} 在 $n = 3$ 维的反例比较, 后者表明存在形如 $\Box u = q(\partial u)$ 的方程(其中 $q$ 是 $\mathbb{R}^4$ 上的二次型), 对 $H^2$ 中的数据不适定. 然而, 若二次型 $q$ 为零形式型, 则几乎可一路推进到临界 WP 指数 $s_c$. 下面两个例子在第一迭代层面验证了这一点.
\begin{example}\label{ex:2.2}
考虑方程$^{6}$
\begin{equation}
\Box u=Q_{0}(u,u),\end{equation}
其中 $u$ 为实值. 同样, 第一迭代的 WP 问题导至证明形如 (24) 的估计, 但由于算子 $Q_{0}$ 的特殊零结构, 奇异因子 $\Delta_{\pm}$ 从核的分母中完全消去. 事实上, $K$ 为
\[
K(\xi,\eta)=\langle\xi\rangle^{-s}+\langle\eta\rangle^{-s},
\]
故由\propref{prop:2.6}, 第一迭代对 $H^{s}$ 中数据 WP, 当 $s > s_{c} = \frac{n}{2}$.\end{example}

\begin{example}\label{ex:2.3}
考虑方程
\begin{equation}
\Box u=Q(u,u),\end{equation}
其中 $u$ 为向量值, $Q(u,u)$ 的第 $I$ 个分量是所有 $i,j,J,K$ 的 $Q_{ij}(u^{J},u^{K})$ 的线性组合. 如上例, 由于 $Q_{ij}$ 的零结构产生消去, 但此情形只消去 $\Delta_{\pm}$ 的半次幂. 事实上, 现在的 $K$ 为
\[
K(\xi,\eta)=\bigl(\langle\xi\rangle^{-s+\frac{1}{2}}+\langle\eta\rangle^{-s+\frac{1}{2}}\bigr)\left(1+\Delta_{\pm}(\xi,\eta)\right)^{-\frac{1}{2}},
\]
故第一迭代对 $H^{s}$ 中数据 WP, 当 $s > \max\left(\frac{n}{2}, \frac{n+3}{4}\right)$.
经显然的修改, 若考虑方程
\[
\Box u=Q(D^{-1}u,u),
\]
则发现第一迭代对 $H^{s}$ 中数据 WP, 当 $s > \max\left(\frac{n-2}{2}, \frac{n-1}{4}\right)$.\end{example}

高阶迭代 一旦第一迭代的分析完成, 真正困难的工作才开始. 为证明迭代 $u_{k+1}$ 满足 (23), 仅用所有先前 $u_j\,(j \leq k)$ 的 (23) 是不够的$^{7}$. 我们需要一个与第一讲所述的 Strichartz 及双线性估计密切相连、强大得多的函数空间.
记号  tempered 分布 $u \in \mathcal{S}'(\mathbb{R}^{n+1})$ 的傅里叶变换记为 $\mathcal{F}u$ 或 $\widehat{u}$, 维数任意. 在频率空间我们用坐标 $(\tau, \xi) = \Xi = (\Xi^0, \ldots, \Xi^n)$, 其中 $\tau \in \mathbb{R}$ 与 $\xi \in \mathbb{R}^n$ 分别对应于时间变量 $t$ 与空间变量 $x$. $\mathbb{R}^{1+n}$ 上的 Lorentz 内积为
\[
\langle\Xi;\widetilde{\Xi}\rangle=\sum_{\mu=0}^{n}\Xi_{\mu}\widetilde{\Xi}^{\mu}=-\Xi^{0}\widetilde{\Xi}^{0}+\sum_{j=1}^{n}\Xi^{j}\widetilde{\Xi}^{j},
\]
而波动算子 $\Box$ 的象征为 $-\langle \Xi;\Xi\rangle = \tau^2 - |\xi|^2$. 我们总以 $|\Xi|$ 表示欧氏范数 $|\Xi|^2 = \tau^2 + |\xi|^2$.
\begin{definition}\label{def:2.8}
设 $\Lambda^{\alpha}$、$\Lambda_{+}^{\alpha}$ 与 $\Lambda_{-}^{\alpha}$ 为如下乘子:
\begin{align}
\widehat{\Lambda^{\alpha}f}(\xi)&=\left(1+|\xi|^{2}\right)^{\alpha/2}\widehat{f}(\xi),\qquad
\widehat{\Lambda_{+}^{\alpha}u}(\Xi)=\left(1+|\Xi|^{2}\right)^{\alpha/2}\widehat{u}(\Xi),\nonumber\\
\widehat{\Lambda_{-}^{\alpha}u}(\Xi)&=\left(1+\frac{\langle\Xi;\Xi\rangle^{2}}{1+|\Xi|^{2}}\right)^{\alpha/2}\widehat{u}(\Xi).\end{align}
\end{definition}

我们定义空间 $H^{s,\theta}(\mathbb{R}^{n+1})$, 它适配于 $\mathbb{R}^{1+n}$ 上的波动算子, 正如 $H^{s}(\mathbb{R}^{n})$ 适配于 $\mathbb{R}^{n}$ 上的 Laplacian.
\begin{definition}\label{def:2.9}
对 $s, \theta \in \mathbb{R}$, 定义
\[
H^{s,\theta}=\left\{u\in\mathcal{S}^{\prime}(\mathbb{R}^{n+1}):\Lambda^{s}\Lambda_{-}^{\theta}u\in L^{2}(\mathbb{R}^{n+1})\right\},
\]
其范数为 $\|u\|_{s,\theta}=\|\Lambda^{s}\Lambda_{-}^{\theta}u\|_{L^{2}}$.\end{definition}

$H^{s,\theta}$ 与具 $H^s$ 中数据的齐次波动方程解空间之间有极简单的联系. 事实上, 每个 $u \in H^{s,\theta}$ 都可写成形式
\[
u(t)=\frac{1}{2\pi}\int_{-\infty}^{\infty}\frac{e^{i t\lambda}u_{\lambda}(t)}{(1+\left|\lambda\right|)^{\theta}}\,d\lambda\qquad(H^{s}\text{ 值积分}),
\]
其中 $\{u_\lambda\}_{\lambda\in\mathbb{R}}$ 是 (1) 的一族单参数解, 数据属于 $H^s$; 即 $\Box u_\lambda = 0$ 且 $(u_\lambda, \partial_t u_\lambda)_{t=0} = (f_\lambda, 0)$, 其中 $\lambda \mapsto f_\lambda$ 属于 $L^2(\mathbb{R}, H^s)$. 而且 $\|u\|_{s,\theta}^2 = \int \|f_\lambda\|^2 d\lambda$.
(28) 的一个重要推论如下:
\begin{theorem*}[原理]
对具 $H^{s}$ 中数据的齐次波动方程解的一个线性或多次线性时空估计, 蕴含对 $\theta > \frac{1}{2}$ 时 $H^{s,\theta}$ 中元素的相应估计.\end{theorem*}

空间 $H^{s,\theta}$ 提供了证明主定理的粗略框架. 事实上它们只适用于证明主定理的 i) 部分. 所有进一步进展都要求对这些空间作实质性精化. 关于 $H^{s,\theta}$ 空间及其修正, 连同第一讲的双线性与 Strichartz 估计与上述原理, 如何用于证明主定理中 $s > s_c$ 的适定性结果, 我们请读者参阅我们的综述 \cite{klai45}. 情形 $s = s_c$ 在 T. Tao 的近期工作之前完全开放, 这将在下一节描述.
\section*{2.10 临界正则性的波动映射}

下面我将讨论 T. Tao 最近得到的关于临界正则性波动映射的突破性结果. 他的结果有力地展示了物理空间中几何思想与傅里叶分析技巧(如仿微分法、Strichartz 与双线性估计)之间的相互作用. 因此它非常契合本讲稿的框架.
他的结果适用于波动映射 $\phi : \mathbb{R}^{n+1} \to S^N$, $S^N$ 为欧氏空间 $\mathbb{R}^{N+1}$ 中的单位球面, 对任意维数 $n \geq 2$. 维数 $n \geq 5$ 的情形容易得多, 从而可专注于 Tao 的主要新思想------规范重整化(gauge renormalization). 事实上我将介绍 Tao 的 $n \geq 5$ 结果由 Rodnianski 与我和 \cite{klai41} 推广到一般目标流形类的一个推广.
设 $\phi: \mathbb{R}^{n+1} \longrightarrow (\mathcal{N}, h)$, 其中 $(\mathcal{N}, h)$ 为维数 $N$ 的黎曼流形. 这里 $\mathbb{R}^{n+1}$ 表示赋予度量 $m = \mathrm{diag}(-1, 1, \ldots, 1)$ 的标准 Minkowski 空间. 记 $\nabla$ 为 $\mathcal{N}$ 的切丛 $T\mathcal{N}$ 上的 Levi-Civita 联络. 在 $\mathcal{N}$ 上的局部坐标 $y^I\,(I = 1, \ldots, N)$ 中, 波动映射方程取熟知形式
\[
\partial^{\alpha}\partial_{\alpha}\phi^{I}+\Gamma_{J K}^{I}\partial_{\beta}\phi^{J}\partial_{\gamma}\phi^{K}m^{\beta\gamma}=0,
\]
其中 $\Gamma_{JK}^{I}$ 为 $\mathcal{N}$ 上 Levi-Civita 联络 $\nabla$ 的 Christoffel 系数, $\phi^{I}\,(I = 1, \ldots, N)$ 为映射 $\phi$ 在 $\mathcal{N}$ 局部坐标中的分量. 令 $e_{a} = e_{a}^{I} \frac{\partial}{\partial y^{I}}$ 为一组正交标架向量场, $\omega^{a} = \omega_{I}^{a} dy^{I}$ 为对应的对偶 1-形式基, $\omega^{a}(e_{b}) = \delta_{b}^{a}$. 由于 $h(e_{a}, e_{b}) = \delta_{ab}$, 推得 $h_{IJ} = \sum_{a} \omega_{I}^{a} \omega_{J}^{a}$. 定义
\[
\phi_{\alpha}^{a}=\omega_{I}^{a}\partial_{\alpha}\phi^{I},
\]
其中 $\phi^{I}$ 为映射 $\phi$ 相对于 $\mathcal{N}$ 上局部坐标 $y^{I}$ 的分量. 显然 $\partial_{\alpha}\phi^{I} = e_{a}^{I}\phi_{\alpha}^{a}$. 对 $\mathcal{N}$ 上函数 $F$ 记 $\partial_{\alpha}F(\phi) = \frac{\partial}{\partial y^{I}}F(\phi)\partial_{\alpha}\phi^{I} = \frac{\partial}{\partial y^{I}}F(\phi)e_{d}^{I}\phi_{\alpha}^{d} = F_{,d}(\phi)\phi_{\alpha}^{d}$, 其中 $F_{,d} = e_{d}(F)$.
我们易验证与波动映射 $\phi$ 相伴的函数 $\phi_{\alpha}^a = \langle \partial_\alpha \phi, e_a \rangle$ 满足$^{8}$ 如下散度--旋度系统:
\[
\partial_{\beta}\phi_{\alpha}^{a}-\partial_{\alpha}\phi_{\beta}^{a}=C_{b c}^{a}\phi_{\alpha}^{b}\phi_{\beta}^{c},
\]
\[
\partial^{\alpha}\phi_{\alpha}^{a}=-\Gamma_{b c}^{a}\phi_{\beta}^{b}\phi_{\gamma}^{c}m^{\beta\gamma},
\]
其中 $C_{bc}^{a}$ 与 $\Gamma_{bc}^{a}$ 分别为标架的结构系数与联络系数,
\[
[e_{b},e_{c}]=C_{b c}^{a}e_{a},
\]
\[
\nabla_{e_{b}}e_{c}=\Gamma_{b c}^{a}e_{a}.
\]
由公式 $[e_b, e_c] = \nabla_{e_b} e_c - \nabla_{e_c} e_b$ 推得
\[
C_{b c}^{a}=\Gamma_{b c}^{a}-\Gamma_{c b}^{a}.
\]
由于标架 $e_a$ 正交, $\Gamma^a_{bc} = \langle \nabla_{e_b} e_c ; e_a \rangle = - \langle e_c ; \nabla_{e_b} e_a \rangle$, 故
\[
\Gamma_{b c}^{a}=-\Gamma_{b a}^{c}.
\]
此外,
\[
\Gamma_{b c}^{a}=\frac{1}{2}\bigl(C_{b c}^{a}+C_{a c}^{b}+C_{a b}^{c}\bigr).
\]

\setcounter{mtthm}{10}% 2.10 为节标题, 下一定义环境显示 2.11

\begin{definition}\label{def:2.11}
若黎曼流形 $\mathcal{N}$ 上存在一组正交标架 $(e_a)_{a}$ 使得结构系数 $C_{bc}^a$ 及其标架导数 $C_{bc,d_1d_2...d_k}^a$ 在 $\mathcal{N}$ 上一致有界, 则称 $\mathcal{N}$ 具有"有界可平行化"(bounded parallelizable)结构.\end{definition}

\begin{remark}\label{rem:2.12}
有界可平行化流形有大量例子. 首先, 在任何李群上我们可构造一组左不变向量场的正交基 $e_a$, 相对它结构常数 $C_{bc}^a$ 为常数$^{9}$. 常负曲率流形 $\mathbf{H}^N\,(N > 2)$ 是有界可平行化的. 而且 $\mathbf{H}^2$(即双曲平面)是一个李群, 见 \cite{klai11} 第 3.1 节的讨论. 此外任何紧黎曼流形都可等距嵌入为一个有界可平行化黎曼流形的全测地子流形, 见 \cite{klai11}.\end{remark}

\begin{proposition}\label{prop:2.13}
设 $\mathcal{N}$ 为黎曼流形, $\phi: \mathbb{R}^{n+1} \to \mathcal{N}$ 为波动映射. 1-形式 $\phi_\alpha^a = \langle \partial_\alpha \phi, e_a \rangle$ 满足方程 (31)、(32) 以及如下波动方程组:
\[
\Box\Phi=-2R_{\mu}\cdot\partial^{\mu}\Phi+E,
\]
其中 $\Phi = (\phi_{\alpha}^{a})$, $R_{\mu} = (R_{b\mu}^{a})_{a,b=1}^{N}$, $R_{b\mu}^{a} = \Gamma_{cb}^{a}\phi_{\mu}^{c}$. 分量 $E = (E_{\alpha}^{a})$ 是 $\Phi = (\phi_{\alpha}^{a})$ 各分量的三次齐次多项式, 其系数只依赖于结构函数 $C_{bc}^{a}$ 及其关于标架的导数 $C_{bc,d}^{a}$.\end{proposition}

\begin{remark}\label{rem:2.14}
必须指出矩阵 $R_{\mu}$ 是反对称的, 即
\[
R_{b\mu}^{a}=-R_{a\mu}^{b}.
\]
这是 (??) 的直接推论. 这表明著名的"Helein 技巧"------在 $N$ 为标准球面的特例下将波动映射方程反对称化, 这一技巧在 Tao 关于标准球面的工作 \cite{klai70}、\cite{klai71} 中起根本作用------事实上源于任何黎曼流形上相对正交标架的联络系数的一般特征.\end{remark}

我们研究波动映射满足初值问题
\[
\phi(0)=\varphi,\qquad\partial_{t}\phi(0)=\psi=\psi_{0}^{a}e_{a}
\]
的演化, 其中 $\varphi$ 为从 $\mathbb{R}^n$ 取值于 $\mathcal{N}$ 的任意光滑映射, $\psi = \psi_0^a e_a$ 为从 $\mathbb{R}^{n+1}$ 到 $T\mathcal{N}$ 的任意光滑映射. 令 $\varphi_i^a = \langle \partial_i \varphi, e_a \rangle$.
\begin{definition}\label{def:2.15}
若所有分量 $\varphi_i^a$、$\psi_i^a$ 属于空间 $\dot{H}^{s-1}(\mathbb{R}^n)$(相应地 $H^{s-1}(\mathbb{R}^n)$), 则称初数据 $\phi[0] = (\varphi, \psi)$ 属于 Sobolev 空间 $\dot{H}^s(\mathbb{R}^n)$(相应地 $H^s(\mathbb{R}^n)$). 我们记
\[
\|\phi[0]\|_{\dot{H}^{s}}=\sum_{a,i}\left(\|\varphi_{i}^{a}\|_{\dot{H}^{s-1}}+\|\psi_{i}^{a}\|_{\dot{H}^{s-1}}\right),
\]
对 $\|\phi[0]\|_{H^{s}}$ 类似.\end{definition}

\begin{theorem}\label{thm:2.16}
设 $\mathcal{N}$ 为装备有界可平行化结构的黎曼流形. 设 $n \geq 5$ 且初数据 $\phi[0] = (\varphi, \psi = \psi_i^a e_a)$ 对某 $\frac{n}{2} < s$ 属于 $H^s$. 我们还作临界小性假设:
\[
\|\phi[0]\|_{\dot{H}^{\frac{n}{2}}}\leq\varepsilon.
\]
则具初数据 $\phi[0]$ 的波动映射 $\phi$ 可在 $H^{s}$ 范数下唯一地整体延拓到时刻 $t$.\end{theorem}

\begin{remark}\label{rem:2.17}
该定理把 \cite{klai70} 的结果从目标流形为标准球面推广到有界可平行化流形. 对维数的限制 $n \geq 5$ 与 \cite{klai70} 中相同; 这使我们只需依赖 Strichartz 估计. 对标准球面的情形, 维数限制由 \cite{klai71} 借助双线性估计$^{10}$、$H^{s,\theta}$ 空间以及 \cite{klai73} 的精化方法移除. 即便有 \cite{klai71}, 将我们的结果推广到二维看来也并非易事$^{11}$.\end{remark}

主定理的证明依赖于 $H^{s}$($s > \frac{n}{2}$)中的局部适定性结果. 我们陈述如下精确结果:
\begin{theorem}\label{thm:2.18}
设初数据 $\phi[0] \in H^s(\mathbb{R}^n)$ 对某 $s \geq s_0 > \frac{n}{2}$. 存在 $T > 0$(只依赖于 $\|\phi[0]\|_{H^{s_0}}$ 的大小)与系统 (31)、(32) 在板 $[0, T] \times \mathbb{R}^n$ 上的唯一解 $\phi$, 满足
\[
\left\|\phi[t]\right\|_{H^{s}}\leq C\|\phi[0]\|_{H^{s}}
\]
对所有 $t \in [0, T]$ 成立, 其中常数 $C$ 只依赖于 $T$、$s$、$s_0 - \frac{n}{2}$ 与 $n$.\end{theorem}

\begin{proof}
我们用 \cite{klai70} 的 Littlewood--Paley 记号. 对函数 $\phi(t, x)$ 记投影
\[
P_k\phi(t, x) = \int e^{ix \cdot \xi} \chi(2^{-k}\xi)\hat\phi(t, \xi) \, d\xi,
\]
其中 $\hat\phi(t, \xi)$ 为 $\phi$ 的空间傅里叶变换, $\chi(\xi) = \eta(\xi) - \eta(2\xi)$, $\eta$ 为支在 $|\xi| \leq 2$ 上、在球 $|\xi| \leq 1$ 上等于 1 的非负光滑鼓包函数. 故 $\chi(\xi)$ 支在 $\{\frac{1}{2} \leq |\xi| \leq 2\}$, 且对 $\xi \neq 0$ 有 $\sum_{k \in \mathbf{Z}} \chi(2^{-k}\xi) = 1$. 又定义 $P_{\leq k} = \sum_{l \leq k} P_l$, 以及对任意区间 $I \subset \mathbf{Z}$ 显然地定义 $P_I$.
记号: 我们将频繁使用记号 $A \lesssim B$ 表示对某不依赖于估计中任何重要参数的常数 $c > 0$ 有 $A \leq cB$.
依 \cite{klai70} 引入记号
\[
\|\Phi\|_{S_{k}}=\sup_{q,r\in\mathcal{A}}2^{k\left(\frac{1}{q}+\frac{n}{r}-1\right)}\left(\|\Phi\|_{L_{t}^{q}L_{x}^{r}}+2^{-k}\|\partial_{t}\Phi\|_{L_{t}^{q}L_{x}^{r}}\right),
\]
其中 $\mathcal{A} = \{(q, r): 2 \leq q, r \leq \infty,\ \frac{1}{q} + \frac{n-1}{2r} \leq \frac{n-1}{4}\}$ 为容许 Strichartz 指数的集合. 回顾:\end{proof}

\begin{proposition}[Strichartz 型估计]\label{prop:2.19}
对任意固定整数 $k$ 与 $\mathbb{R} \times \mathbb{R}^n$ 上函数 $\phi(t, x)$, 若其 $\hat{\phi}(t, \xi)$ 的支包含于二进区域 $2^{k-1} \leq |\xi| \leq 2^{k+1}$, 则有估计
\[
\|\phi\|_{S_{k}}\lesssim\|\phi(0)\|_{H^{\frac{n-2}{2}}}+\|\partial_{t}\phi(0)\|_{H^{\frac{n-4}{2}}}+2^{k\frac{n-4}{2}}\|\Box\phi\|_{L_{t}^{1}L_{x}^{2}}.
\]
\end{proposition}

下面我们回顾 \cite{klai70} 中给出的频率包络定义.
\begin{definition}\label{def:2.20}
频率包络是一个 $l^2$ 序列 $c = (c_k)_{k \in \mathbf{Z}}$ 满足
\[
c_{k}\lesssim2^{\sigma|k-k^{\prime}|}c_{k^{\prime}}
\]
对所有 $k, k' \in \mathbf{Z}$ 成立. 这里 $\sigma$ 为固定的正常数; 如 \cite{klai70} 取 $0 < \sigma < \frac{1}{2}$. 此外我们需要 $0 < \sigma < \frac{n-4}{4}$ 与 $0 < \sigma < \frac{n-3}{4(n-1)}$.\end{definition}

若 $\mathbb{R}^n$ 上函数 $f$ 的 $\hat{H}^s$ 范数位于包络 $c$ 之下, 是指对所有 $k \in \mathbf{Z}$ 有 $\|P_k f\|_{\dot{H}^s} \leq c_k$. 记为 $f \ll_s c$, 无歧义时简记 $f \ll c$. 易验证若 $\|f\|_{\dot{H}^s} \leq \varepsilon$, 则存在包络 $c \in l^2$ 使 $\|c\|_{l^2} \lesssim \varepsilon$ 且 $f \ll_s c$. 事实上我们可取 $c_k = \sum_{k' \in \mathbf{Z}} 2^{-\sigma |k - k'|} \|P_k\|_{\dot{H}^s}$.
\begin{definition}\label{def:2.21}
固定 $0 < \sigma < \min(\frac{1}{2}, \frac{n-4}{4}, \frac{n-3}{4(n-1)})$ 与频率包络 $c$. 若相对我们的标架 $e_a$, 对所有 $k \in \mathbf{Z}$ 有
\[
\left\|P_{k}\phi[0]\right\|_{\dot{H}^{\frac{n}{2}}}\leq c_{k},
\]
则称初数据 $\phi[0] = \left( \phi(0) = \varphi, \partial_t \phi(0) = \psi = \psi_0^a e_a \right)$ 位于 $c$ 之下. 简记为 $\phi[0] \ll c$.\end{definition}

我们的主定理的证明可轻易化为如下:
\begin{proposition}[主命题]\label{prop:2.22}
设 $c$ 为频率包络$^{12}$, 满足 $\|c\|_{l^2} \leq \varepsilon$, $0 < T < \infty$, 且 $\Phi = (\phi_\alpha^a)$ 满足方程 (31)、(32), 从而也满足 (34). 依\defref{def:2.21}, 设初数据满足小性条件 $\phi[0] \ll c$. 再设自举(bootstrap)假设
\[
\left\|P_{k}\Phi\right\|_{S_{k}([0,T]\times\mathbb{R}^{n})}\leq2C c_{k}
\]
对所有 $k \in \mathbf{Z}$ 成立. 则事实上对充分小的 $\varepsilon$ 与所有 $k \in \mathbf{Z}$,
\[
\|P_{k}\Phi\|_{S_{k}([0,T]\times\mathbb{R}^{n})}\leq C c_{k}.
\]
\end{proposition}

\begin{remark}\label{rem:2.23}
由于我们的方程与小性条件 $\phi[0] \ll c$ 的尺度不变性, 只需对 $k = 0$ 证明 (40). 令 $\Psi = P_0\Phi$. 我们需证明
\[
\|\Psi\|_{S_{0}([0,T]\times\mathbb{R}^{n})}\leq C c_{0}.
\]
\end{remark}

\begin{remark}\label{rem:2.24}
为证明 (41), 我们想把\propref{prop:2.19} 应用于对 (34) 作投影 $P_{0}$ 所得的方程, 即
\[
\Box\Psi=P_{0}(R_{\mu}\cdot\partial^{\mu}\Phi+E).
\]
然而 Strichartz 不等式的直接应用不会奏效. 事实上依\propref{prop:2.19},
\[
\|\Psi\|_{S_{0}}\leq c_{0}+\|P_{0}(R_{\mu}\cdot\partial^{\mu}\Phi+E)\|_{L_{t}^{1}L_{x}^{2}}.
\]
三次项 $E$ 不成问题, 困难出现在我们试图更精细地估计 $P_0(R_\mu \cdot \partial^\mu \Phi)$------更精确地, 出现在 $R$ 的低频与同 $\Psi$ 频率相当的 $\Phi$ 频率相互作用的部分. 最危险的项形如 $\tilde{R} \cdot \partial \Psi$, 其中 $\tilde{R} = P_{\leq -10} R$. 为相对于可用的 Strichartz 范数估计 $\|\tilde{R} \cdot \partial \Psi\|_{L_t^1 L_x^2}$, 我们被迫取 $\Psi$ 在能量范数 $L_t^\infty L_x^2$ 中. 这留下 $\tilde{R}$ 在范数 $L_t^1 L_x^\infty$ 中, 而对此我们没有 Strichartz 估计.\end{remark}

\begin{definition}\label{def:2.25}
下面我们用 Tao 的约定, 称 $[0, T] \times \mathbb{R}^n$ 上任一函数或矩阵值函数 $F$ 为可接受误差(acceptable error), 如果
\[
\|F\|_{L_{t}^{1}L_{x}^{2}([0,T]\times\mathbb{R}^{n})}\leq C^{3}\varepsilon c_{0}.
\]
\end{definition}

\begin{proposition}\label{prop:2.26}
矩阵值函数 $P_{0}\Phi = \Psi$ 满足方程
\[
\Box\Psi=-2\tilde{R}_{\mu}\cdot\partial^{\mu}\Psi+\text{error},
\]
其中 $\tilde{R}_{\mu} = P_{\leq -10} R_{\mu} = \Gamma \cdot \tilde{\Phi}_{\mu}$, $\tilde{\Phi}_{\alpha} = P_{\leq -10} \Phi_{\alpha}$. 这里"error"指 (42) 意义下的可接受误差项.\end{proposition}

\begin{proposition}\label{prop:2.27}
矩阵值函数 $\Psi$ 满足形如
\[
\Box\Psi=-2\partial_{\mu}\tilde{\Delta}\cdot\partial^{\mu}\Psi+\text{error}
\]
的方程, 其中势 $\Delta$ 满足如下性质:
\begin{enumerate}[label=\roman*.)]
\item $N \times N$ 矩阵 $\tilde{\Delta}$ 反对称, 即 $\tilde{\Delta}^t = -\tilde{\Delta}$. $\tilde{\Delta}$ 各分量的空间傅里叶变换支在 $|\xi| \leq 2^{-10}$.
\item 对任意 $\tilde{\Delta}_{k} = P_{k}\tilde{\Delta}$ 成立估计:
\[
\|\tilde{\Delta}_{k}\|_{S_{k}}\lesssim 2^{-k}C c_{k},\qquad
\|\partial\tilde{\Delta}_{k}\|_{S_{k}}\lesssim C c_{k},\qquad
\|\Box\tilde{\Delta}_{k}\|_{S_{k}}\lesssim2^{k}C c_{k}.
\]
\item 令 $\bar{R}_{\mu} = \tilde{R}_{\mu} - \partial_{\mu} \tilde{\Delta}$. 对所有 $P_{k} \bar{R}$ 成立:
\[
\|P_{k}\bar{R}\|_{L_{t}^{1}L_{x}^{\infty}}\lesssim C^{2}c_{k}^{2},\qquad
\|P_{k}\bar{R}\|_{L_{t}^{\infty}L_{x}^{\infty}}\lesssim 2^{k}C^{2}c_{k}^{2}.
\]
\end{enumerate}
\end{proposition}

设 $M$ 为依赖于 $T$ 的大整数, 将在下面选取. 定义实 $N \times N$ 矩阵值函数 $U$ 为
\[
U=I+\sum_{-M<k\leq-10}U_{k},
\]
其中 $U_{k}$ 归纳定义如下:
\[
U_{k}=0\quad\text{对所有}\ k<-M;\qquad U_{-M}=I;\qquad
U_{k}=\tilde{\Delta}_{k}\cdot U_{<k}\quad\text{对所有}\ -M<k\leq-10,
\]
其中 $U_{<k} = \sum_{k' < k} U_{k'}$. 由于矩阵 $\tilde{\Delta}_k = P_k\tilde{\Delta}$ 反对称, 得到恒等式
\[
U_{k}^{t}\cdot U_{<k}+U_{<k}^{t}\cdot U_{k}=0,
\]
从而
\[
U_{<k}^{t}\cdot U_{<k}-I=\sum_{k^{\prime}<k}U_{k^{\prime}}^{t}\cdot U_{k^{\prime}}.
\]
利用此恒等式可归纳证明
\[
\|U_{<k}\|_{L_{t}^{\infty}L_{x}^{\infty}}\leq 2,\qquad
\|U_{k}\|_{L_{t}^{\infty}L_{x}^{\infty}}\lesssim C c_{k}\quad (k>-M),
\]
以及
\[
\|U_{k}\|_{L_{t}^{2}L_{x}^{\infty}}\lesssim C2^{-k/2}c_{k}\quad (k>-M).
\]
此外,
\[
\|\partial U_{<k}\|_{L_{t}^{\infty}L_{x}^{\infty}}\lesssim 2^{k}C^{2}c_{k},\qquad
\|\partial U_{k}\|_{L_{t}^{\infty}L_{x}^{\infty}}\lesssim 2^{k}C^{2}c_{k},
\]
以及
\[
\|\partial U_{<k}\|_{L^{2}_{t} L^\infty_{x}}\lesssim 2^{k/2}C c_{k},\qquad
\|\partial U_{k}\|_{L^{2}_{t} L^\infty_{x}}\lesssim 2^{k/2}C c_{k}.
\]

\begin{proposition}\label{prop:2.28}
设 $\varepsilon$ 充分小(依赖于 $C$), 且 $M$ 充分大(依赖于 $T, C, \varepsilon$). 则矩阵 $U$ 满足如下性质:
\begin{enumerate}[label=\roman*.)]
\item 近似正交性:
\[
\|U^{t}U-I\|_{L_{t}^{\infty}L_{x}^{\infty}},\quad\|\partial(U^{t}U-I)\|_{L_{t}^{\infty}L_{x}^{\infty}}\lesssim C^{2}\varepsilon.
\]
特别地, 对小 $\varepsilon$, $U$ 可逆且
\[
\|U\|_{L_{t}^{\infty}L_{x}^{\infty}},\quad\|U^{-1}\|_{L_{t}^{\infty}L_{x}^{\infty}}\lesssim1.
\]
\item 近似规范条件:
\[
\|\partial_{\mu}U-\partial_{\mu}\tilde{\Delta}\cdot U\|_{L_{t}^{1}L_{x}^{\infty}}\lesssim C^{2}\varepsilon.
\]
\item 此外,
\[
\|\partial U\|_{L_{t}^{\infty}L_{x}^{\infty}},\quad\|\partial U\|_{L_{t}^{1}L_{x}^{\infty}}\lesssim C^{2}\varepsilon,\qquad
\|\Box U\|_{L_{t}^{2}L_{x}^{n-1}}\lesssim C^{2}\varepsilon.
\]
\end{enumerate}
\end{proposition}

依 \cite{klai70} 我们现在可以进行规范变换
\[
\Psi=U\cdot W.
\]
$W$ 满足方程
\[
\Box W=-2U^{-1}(\partial_{\mu}U-\partial_{\mu}\tilde{\Delta}\cdot U)\partial^{\mu}W
-2U^{-1}\partial_{\mu}\tilde{\Delta}\cdot(\partial^{\mu}U)U^{-1}\Psi-U^{-1}(\Box U)U^{-1}\Psi+\text{error}.
\]
由\propref{prop:2.28} 推得:
\begin{proposition}\label{prop:2.29}
矩阵值函数 $W$ 满足形如
\[
\Box W=\text{error}
\]
的方程. 因此, 若 $\varepsilon$ 充分小,
\[
\|\Psi\|_{S_{0}}\lesssim\|W\|_{S_{0}}\leq\|\Psi[0]\|_{H^{\frac{n-2}{2}}}+C C^{3}\varepsilon c_{0}\leq C c_{0}.
\]
这正是 (41), 从而结束主命题的证明.\end{proposition}
